\section{Protocol details}
\label{app:protocol}

\subsection{Berka: the balance field, adapter decisions and accounts}
\emph{The balance field.} Each record carries the bank's balance after it.
We do not use that field as an input or as truth. Instead, the end-of-day
balance is the account's opening balance (the first record's balance minus
its amount) plus the running sum of amounts, carried forward over days
without records. The balance field serves to validate that running sum.
The export stores amounts and balances as whole crowns, rounded
separately, so the two drift apart slightly. Among the held-out accounts
we keep, the largest difference on any day is \R{ho/data/max_drift_czk}
crowns, and the median account's largest difference is
\R{ho/data/median_account_max_drift_czk}. In some accounts the reported
balance resets mid-history.

Two adapter decisions were made on development data and committed before
any development result was read:
\begin{enumerate}[label=(\roman*)]
\item \emph{Calendar.} Berka bookings fall evenly across weekdays,
  including dates a Czech holiday calendar would skip. With that calendar,
  no income schedule could be detected. The engine gained a pluggable
  holiday calendar, and Berka uses an every-day calendar.
\item \emph{Balance drift.} An account is excluded when, on any day, its
  running sum differs by more than \R{c/max-drift}~crowns from every
  balance the bank reported that day. The rule uses each account's full
  history, which an application could not do at a forecast origin; it is a
  data-quality selection made once, not a forecast input.
\end{enumerate}
The approval of rule (ii) described it, wrongly, as a threshold of
\R{c/drift-high}~crowns. The code and the plan, committed together before any
results, use \R{c/max-drift}. Table~\ref{tab:drift} gives the exclusions at
both thresholds. The held-out set is identical under either. On
development, the applied rule excluded one account that a \R{c/drift-high}-crown
rule would have kept.

\begin{table}[t]
\centering\footnotesize
\begin{tabular}{l rrrr}
\toprule
Split & accounts & drift $>$ \R{c/drift-low} & drift $>$ \R{c/max-drift} (applied) & drift $>$ \R{c/drift-high}\\
\midrule
Development & \R{au/drift/development/accounts} & \R{au/drift/development/drift_over_50} & \R{au/drift/development/drift_over_100} & \R{au/drift/development/drift_over_1000}\\
Held-out & \R{au/drift/held-out/accounts} & \R{au/drift/held-out/drift_over_50} & \R{au/drift/held-out/drift_over_100} & \R{au/drift/held-out/drift_over_1000}\\
\bottomrule
\end{tabular}
\caption{Accounts excluded for balance drift (crowns) at three thresholds,
recounted from the raw data with the code's rule.}
\label{tab:drift}
\end{table}

\emph{Accounts used.} Of the \R{ho/data/accounts_in_split} held-out
accounts, \R{ho/data/accounts_excluded_balance_inconsistent} are excluded
for drift. Forecasts are attempted for \R{ho/data/accounts_run}, and
\R{ho/data/accounts_with_eligible_origin} of them have at least one
eligible origin (\R{ho/data/eligible_origins} origins). On development the
same steps give \R{dev/data/accounts_in_split},
\R{dev/data/accounts_excluded_balance_inconsistent},
\R{dev/data/accounts_run}, \R{dev/data/accounts_with_eligible_origin} and
\R{dev/data/eligible_origins}.


\subsection{MoneyData}
This is one UK current account, from
\R{md/data/first_day} to \R{md/data/last_day}
\citep{firat2023moneyvis,laramee2026moneydata}, with \R{md/data/records}
records, used as a descriptive case study only. In the distributed
spreadsheet, dates whose day was at most twelve had been read with day and
month swapped. For example, the file's \texttt{\R{d/md-date-raw}} is
\texttt{\R{d/md-date-repaired}}. After the repair, the dates increase with
the record numbers and the balance column chains exactly. Records are
tagged at each origin by the application's own tagger. The default US
calendar is used, because the engine has no UK calendar.


\subsection{Origins, inputs and timing}
Origins fall every
\R{c/origin-step} days from day \R{c/warmup} of an account's history, while
origin plus \hC\ days stays inside the data. An origin needs at least
\R{c/min-records} records in the \R{c/warmup} days before it.
\begin{itemize}
\item \textbf{What a forecast uses.} The records dated on or before the
  origin, including all records of the origin day. Records are ordered by
  date and record id.
\item \textbf{Today's balance.} The end-of-day balance of the origin, the
  last balance carried forward if the origin has no records. The truth is
  the end-of-day balance \hA, \hB\ or \hC\ days later.
\item \textbf{Outcomes known at the origin.} B-SA+'s past errors and V1's
  calibration outcomes use only forecast dates $g$ with $g + h$ on or
  before the origin.
\item \textbf{Prefix test.} For \R{au/prefix/origins} randomly drawn
  development origins (\R{au/prefix/accounts} accounts), we cut the input
  file at the origin and recomputed every per-account method. All
  \R{au/prefix/truncated_equals_full} forecasts equal those made from the
  full file, and all \R{au/prefix/row_equals_frozen} scored rows are
  byte-identical to the frozen development results.
\item \textbf{The non-causal exceptions.} The drift exclusion above. The
  V3/V4 priors and B-GLOBAL's training use other accounts' data from all
  dates, so they carry calendar-future information for early origins.
  Neither affects V0, V1 or V2.
\item \textbf{Seeds.} The forecast seed is a fixed hash of (account,
  origin).
\end{itemize}


\subsection{All baselines}
\begin{itemize}
\item \textbf{B-30} projects the \R{c/b30-window}-day mean daily change,
  with a Gaussian band.
\item \textbf{B-LIN} extends a least-squares line through the last
  \R{c/blin-window} days.
\item \textbf{B-SA} is schedule-aware and deterministic: the same detected
  schedule as the model, the last paycheck and bill amounts, and the
  \R{c/bsa-window}-day mean of other flows.
\item \emph{B-SA+}, the planned reference, adds to B-SA the empirical
  distribution of its own past errors.
\item \textbf{B-SEAS} is a monthly-seasonal naive forecast with its own
  past errors.
\item \emph{B-EQ} adds the account's own past $h$-day balance changes to
  today's balance.
\item \emph{B-GLOBAL} is LightGBM quantile regression
  \citep{ke2017lightgbm} on lagged balance and flow features, with fixed
  hyperparameters (\R{c/lgbm-rounds} rounds). It is cross-fitted over
  development accounts and, for the held-out run, trained on all of them.
\end{itemize}
B-GLOBAL pools data across accounts, which the local-first design forbids.
It is a reference point, not a like-for-like competitor; B-EQ, which uses
only the account's own history, is.


\subsection{Scoring details}
\label{sec:scoring}
\paragraph{Forecast types.} Each method's forecast for a day is scored as
one of five kinds:
\begin{itemize}
\item \textbf{Monte Carlo cones} (V0, V2, V3) give type-7 quantiles
  \citep{hyndman1996} of \R{c/paths} paths, on a grid of \R{c/crps-grid} midpoint levels
  plus the interval ends.
\item \textbf{Quantile functions} (V1, V4, B-GLOBAL) have \R{c/qf-levels} levels from
  \R{c/qf-lo}\% to \R{c/qf-hi}\%. They are linear between levels and flat beyond them, and
  crossing quantiles are repaired by sorting.
\item \textbf{Empirical samples} (B-SA+, B-SEAS, B-EQ) get the exact CRPS of
  the sample. Their interval ends are order statistics.
\item \textbf{Gaussian} forecasts (B-30, B-LIN) get the closed-form CRPS.
\item \textbf{B-SA} is a point, so its coverage is undefined.
\end{itemize}
A quantile function's CRPS is the pinball-loss integral, taken by the
midpoint rule over \R{c/crps-grid} levels. The flat tails of V1, V4 and B-GLOBAL
therefore understate their tail spread relative to the cones. PIT inverts
the forecast's distribution function, at the middle of any jump or flat
stretch.

\paragraph{V1's quantile function.} Each of V1's central levels has its own
recursion, so the adjusted levels need not increase with the nominal
level. The resulting values are sorted into one monotone function. Where
two levels crossed, their intervals are exchanged.

\paragraph{Coverage, skill and intervals.}
\begin{itemize}
\item \textbf{Coverage} of the central $\ell$ interval counts
  $q_{(1-\ell)/2} \le y \le q_{(1+\ell)/2}$.
\item \textbf{CRPS skill} of $m$ against a reference is aggregated two
  ways, over the forecasts both made:
  \begin{itemize}
  \item the \emph{account mean},
    $\operatorname{mean}_a \bigl(1 - \sum_{a} \mathrm{CRPS}_m / \sum_{a}
    \mathrm{CRPS}_{\text{ref}}\bigr)$, which weights accounts equally and
    was the planned estimand;
  \item the \emph{pooled} ratio
    $1 - \sum \mathrm{CRPS}_m / \sum \mathrm{CRPS}_{\text{ref}}$ over all
    origins, which weights origins equally and, unlike the account mean,
    composes across references and strata.
  \end{itemize}
\item \textbf{Intervals} are \R{c/ci-level}\% intervals from
  \R{c/bootstrap} bootstrap resamples of whole accounts (seeds, on
  synthetic data).
\end{itemize}
All methods share one origin set on Berka: every method has the same
\R{ho/var/V2/30/forecasts} held-out forecasts in
\R{ho/var/V2/30/accounts} accounts at each horizon.


\subsection{The timeline}
\begin{table}[h]
\centering\scriptsize
\tab{timeline}
\caption{Protocol timeline. Commit times are from git; ``first on GitHub''
is the time of the first push containing the commit, from GitHub's
repository activity record. The held-out runs carry the times their guard
logged.}
\label{tab:timeline}
\end{table}


\subsection{Deviations and post-hoc analyses}
\begin{enumerate}[label=(\arabic*)]
\item The every-day calendar and the balance-drift exclusion for Berka,
  approved on development data before the held-out run. The approval's
  description of the drift threshold was wrong (Table~\ref{tab:drift}).
\item The revision itself. The plan anticipated revision only as a
  possibility; the rule was written after V0 failed, but before any
  revision code or revision result existed.
\item Held-out run~1 omitted the planned ablations through a harness gap:
  the ablation sample was drawn from the development split only. They were
  run as a separately logged second run, \emph{after} the main held-out
  results were known (\S\ref{sec:ablations}).
\item Computed after the held-out results were known, from the stored
  forecasts and without new forecasts:
  \begin{itemize}
  \item the paired comparisons of V2 with each baseline
    (Table~\ref{tab:pairwise});
  \item the pooled skills and mean CRPS reported alongside them;
  \item the one-sided misses (Table~\ref{tab:widths});
  \item the per-cell verdicts for A2 (Table~\ref{tab:a2});
  \item the estimand audit;
  \item the ablations' paired differences and their alarm strata.
  \end{itemize}
\item The synthetic panel was defined after the held-out results:
  \R{c/panel-seeds} seeds of \R{c/panel-days} days, origins every
  \R{c/panel-step} days from day \R{c/panel-first-origin}. The A1 check
  reuses the exact protocol of the engine's earlier calibration test and
  reproduces its numbers.
\item The MoneyData adapter (the date repair and the calendar) was set
  before MoneyData was scored, but after the Berka results.
\item The bit-identity fix (\S\ref{sec:identity}) came after the freeze. The
  frozen results were not regenerated.
\item The analyses of \S\ref{sec:explanation} and \S\ref{sec:overdraft},
  the V2 ablations and B-GLOBAL without calendar look-ahead were specified
  after the held-out results were known, in the repository's
  \texttt{paper/REVISION.md}. They were developed on the development split
  and run on held-out accounts in two separately logged runs. Their
  results are exploratory.
\item Adding B-EQ rows to the development strata table, before the
  held-out run, shifted other cells' bootstrap bounds slightly, because one
  random stream spans all cells. Point estimates did not change.
\end{enumerate}
